\scanpage{40}
\mspage{38}
% archive sheet number: 39
($V_\alpha$ régulier, \deleted{$\operatorname{div}(t_{V_\alpha})$ rég.} $V_\alpha/tV_\alpha$ régulier) sont \uline{cofinaux} parmi les $V_\alpha$. De plus, tout $V_\alpha$ a un corps résiduel de type fini. On est donc \unsure{sous|dans} les conditions de 8.1.

\uline{Remarque} 8.3. \deleted{\unsure{Cette dé}} La démonstration du cas b) \deleted{serait} serait valable sans restriction sur $V$, une fois qu'on disposerait de la résolution des singularités \inserted{sous la forme} de Hironaka \deleted{\unsure{sur le}} (avec diviseurs à croisements normaux $\ldots$) pour les schémas de \uline{type fini sur $\unsure{\mathbb Z|Z}$}; \inserted{cela impliquerait que l'hyp. de 8.1. sur $V$ est \uline{\unsure{tjs|tjrs}} satisfaite.} Elle marcherait de même si $V$ est de car. $p>0$, si on dispose de la résolution des singularités à la Hironaka pour les schémas de t.f. sur $\mathbb F_p$. Notons qu'on n'a besoin ici \deleted{\unsure{de la de Hi}} que de ``l'\uline{uniformisation locale}''.

8.4. \uline{Démonstration de} 8.1. On \deleted{\unsure{aura}} \inserted{peut supposer les $V_\alpha$ henséliens,} quitte à les remplacer par leurs hensélisés. On aura
\[
 K=\underset{\longrightarrow}{\lim}\,\underbrace{(V_\alpha)_t}_{K_\alpha},
\]
donc pour $\alpha$ assez grand, $X_K$ provient d'un schéma \deleted{de type f.} $X_{\deleted{\unsure{K}}\alpha}$ sur \deleted{\unsure{Spec}} $(V_\alpha)_t$, i.e. sur l'ouvert $U_\alpha=S_{\alpha t}$ de $S_\alpha=\operatorname{Spec}V_\alpha$; \deleted{\unsure{prenons}} soit $\eta_\alpha$ le pt générique de $S_\alpha$ i.e. de $S_{\alpha t}$, c'est l'image de $\eta$.
\deleted{Comme $f_\alpha:X_{K_\alpha}\to U_\alpha$ est \unsure{Choisissons un ouvert},}
Comme $f_\alpha$ est propre, on sait (SGA 5 \uline{VI}) que $\unsure{E_{U_\alpha}|\mathscr E_{U_\alpha}}=R^if_{\alpha *}(\mathbb Z_\ell)$ sur $U_\alpha$ est constructible, donc \deleted{$\exists$ \unsure{un}} voisinage ouvert $U'_\alpha$ de $\eta_\alpha$ sur lequel $E_\alpha$ est \unsure{constant tordu}. Augmentant au besoin $\alpha$, \inserted{comme $g^{-1}(U_\alpha)=\unsure{g|j}^{-1}(\unsure{V_\alpha|U'_\alpha})=\{\eta\}$,} on
\diagram{\begin{tikzcd}[ampersand replacement=\&]
S \arrow[d,"g"'] \& \eta_\alpha \\
S_\alpha\supset U_\alpha \&
\end{tikzcd}}
